\documentclass[aspectratio=169]{beamer}
\usetheme{default}
\setbeamertemplate{navigation symbols}{}
\usepackage{amsmath,amssymb,amscd}
\title{Reading seminar}
\subtitle{Session 1}
\date{Deformation viewpoint}
\begin{document}
\begin{frame}\titlepage\end{frame}
\begin{frame}{1. Definition of a deformation}
Following Huybrechts, a deformation of a compact complex
manifold $X$ is a proper smooth holomorphic family
\[
\begin{CD}
X @>>> \mathcal X\\
@VVV @VV{\pi}V\\
\{0\} @>>> B
\end{CD}
\]
This is a Cartesian square: $X\simeq\mathcal X\times_B\{0\}$.
The fibres $X_t$ vary the complex structure.

A tangent vector $v\in T_0B$ is an infinitesimal direction.
The Kodaira--Spencer map turns it into a class on $X$.
\end{frame}

\begin{frame}{2. The Hilbert scheme: from functor to space}
\small
Fix a projective variety $Y$ and a Hilbert polynomial $P$.
For each scheme $U$, consider flat families of subschemes
\[
Z\subset Y\times U,\qquad Z\longrightarrow U,
\]
whose fibres have Hilbert polynomial $P$. Pullback of families
along $U\to U'$ defines a contravariant functor
\[
\operatorname{Hilb}_Y^P:(\mathbf{Sch})^{\mathrm{op}}\to\mathbf{Sets}.
\]
Grothendieck's representability theorem constructs
$\operatorname{Hilb}_Y^P$ with a natural bijection
\[
\operatorname{Hom}(U,\operatorname{Hilb}_Y^P)
\ \simeq\ \operatorname{Hilb}_Y^P(U).
\]
The identity map corresponds to the \emph{universal family}
$\mathcal Z\to\operatorname{Hilb}_Y^P$; every $Z\to U$ is
its pullback under a unique classifying map
$U\to\operatorname{Hilb}_Y^P$.
\end{frame}

\begin{frame}{4. The abstract deformation functor $\operatorname{Def}_X$}
\small
Now fix a compact complex manifold $X$, without
choosing an ambient projective variety.
Let $\mathbf{Art}_{\mathbb C}$ be the category of local
Artinian $\mathbb C$-algebras with residue field $\mathbb C$.
Define
\[
\operatorname{Def}_X:\mathbf{Art}_{\mathbb C}\to\mathbf{Sets},
\quad
A\longmapsto
\left\{
\begin{array}{c}
\text{flat deformations }\mathcal X_A\to\operatorname{Spec}A\\
\text{with special fibre identified with }X
\end{array}\right\}\big/\cong.
\]
A map $A\to A'$ sends a deformation to its base change
over $\operatorname{Spec}A'$.

\medskip
Unlike the Hilbert functor, this parametrizes
\emph{abstract} deformations of $X$, not embedded
subschemes of a fixed $Y$. Its local moduli object is
the Kuranishi germ $\operatorname{Def}(X)$.
\end{frame}

\begin{frame}{3. Analytic versus formal deformation spaces}
\small
\textbf{Huybrechts' $\operatorname{Def}(X)$} is the analytic
Kuranishi \emph{germ} at $0$, carrying a semiuniversal
holomorphic family. For IHS $X$, BTT makes it smooth:
\[
(\operatorname{Def}(X),0)\simeq(\mathbb C^N,0),
\qquad N=h^{1,1}(X).
\]
Its \emph{formal completion} is a different object:
\[
\widehat{\operatorname{Def}}(X)=\operatorname{Spf}R,\quad
R=\widehat{\mathcal O}_{\operatorname{Def}(X),0}
=\varprojlim_n\mathcal O_{\operatorname{Def}(X),0}/\mathfrak m^{n+1}
\simeq\mathbb C[[t_1,\ldots,t_N]].
\]
$\operatorname{Spf}R$ is the formal spectrum: the compatible
system of infinitesimal neighborhoods
$\operatorname{Spec}(R/\mathfrak m^{n+1})$.
For IHS manifolds it prorepresents $\operatorname{Def}_X$
(deformations over Artinian bases correspond uniquely to maps
from $R$). By contrast, $\operatorname{Hilb}(Y)$ globally
represents the embedded Hilbert functor.
\end{frame}

\begin{frame}{5. The relative tangent sequence}
Because $\pi$ is a submersion, its differential is surjective.
The kernel consists of vectors tangent to the fibres:
\[
0\longrightarrow T_{\mathcal X/B}
\longrightarrow T_{\mathcal X}
\xrightarrow{d\pi}\pi^*T_B\longrightarrow0.
\]
Restrict to the central fibre. Since
$T_{\mathcal X/B}|_X=T_X$ and
$\pi^*T_B|_X=T_0B\otimes\mathcal O_X$, we obtain
\[
\boxed{0\longrightarrow T_X\longrightarrow
T_{\mathcal X}|_X\longrightarrow
T_0B\otimes\mathcal O_X\longrightarrow0.}
\]
\end{frame}

\begin{frame}{6. The connecting homomorphism trick}
Apply the long exact sequence of sheaf cohomology:
\[
H^0(X,T_{\mathcal X}|_X)\longrightarrow
\underbrace{H^0(X,T_0B\otimes\mathcal O_X)}_{=\,T_0B}
\xrightarrow{\ \delta\ }H^1(X,T_X)\longrightarrow\cdots
\]
Here $H^0(X,\mathcal O_X)=\mathbb C$ for compact connected $X$.
The connecting map is precisely Kodaira--Spencer:
\[
\boxed{\rho=\delta:T_0B\longrightarrow H^1(X,T_X).}
\]
A local lift of $v$ to $T_{\mathcal X}|_X$ need not glue
globally. Differences of local lifts are tangent to $X$
and define a Cech $1$-cocycle representing $\rho(v)$.
\end{frame}
\begin{frame}{7. The IHS computation}
Dolbeault gives $H^1(X,T_X)\simeq
H_{\bar\partial}^{0,1}(X,T^{1,0}X)$.
For an IHS manifold, contraction with the nondegenerate
holomorphic symplectic form $\sigma$ is an isomorphism
\[
T_X\xrightarrow{\ \iota_{(-)}\sigma\ }\Omega_X^1.
\]
Hence, by Dolbeault and the Kahler Hodge decomposition,
\[
\boxed{T_0\operatorname{Def}(X)
\simeq H^1(X,T_X)\simeq H^1(X,\Omega_X^1)
\simeq H^{1,1}(X).}
\]
The Bogomolov--Tian--Todorov theorem ensures that
$\operatorname{Def}(X)$ is smooth.
\end{frame}
\begin{frame}{8. Step zero: identify the cohomology lattices}
\small
Let $\pi:\mathcal X\to B$ be a smooth deformation of an IHS
manifold $X=X_0$. The groups $H^2(X_t,\mathbb Z)$ form a
local system $R^2\pi_*\mathbb Z$.
Over a simply connected small base, parallel transport identifies
them compatibly with a fixed Beauville--Bogomolov lattice
$(\Lambda,q)$:
\[
\varphi_t:(H^2(X_t,\mathbb Z),q_{X_t})
\xrightarrow{\sim}(\Lambda,q).
\]
A \emph{marking} is such an isometry for one fibre.
The lattice stays fixed, but the Hodge decomposition varies.
\end{frame}

\begin{frame}{9. Periods: moving lines in a fixed vector space}
\small
Each fibre has a one-dimensional space of holomorphic
symplectic forms:
\[
H^{2,0}(X_t)=\mathbb C\sigma_t.
\]
Only the line matters: $\sigma_t$ and $a\sigma_t$ ($a\ne0$)
give the same period. Thus the \emph{period map} is
\[
\mathcal P:t\longmapsto
[\varphi_t(\sigma_t)]\in\mathbb P(\Lambda_{\mathbb C}).
\]
Its image lies in the \emph{period domain}
\[
\Omega_\Lambda=\{[z]:q(z)=0,\ q(z,\bar z)>0\}.
\]
This is a smooth, Hausdorff, noncompact complex manifold
of dimension $b_2(X)-2$.
\end{frame}

\begin{frame}{10. Period domain versus marked moduli}
\small
The marked moduli space remembers pairs $(X,\varphi)$;
the period domain remembers their Hodge-theoretic periods:
\[
\mathcal M_\Lambda\ \xrightarrow{\ \mathcal P\ }\ \Omega_\Lambda,
\qquad (X,\varphi)\longmapsto[\varphi(H^{2,0}(X))].
\]
\emph{Local Torelli:} $\mathcal P$ is a local biholomorphism,
so nearby periods describe nearby marked deformations.

\medskip
\emph{Globally:} marked moduli can be non-Hausdorff, even
though $\Omega_\Lambda$ is Hausdorff. Distinct birational
models may have the same period; global Torelli relates
periods and birational geometry, with marking/monodromy
compatibility essential.
\end{frame}

\begin{frame}{11. Deforming a line bundle (Huybrechts, 1.14)}
\small
Fix $L$ on $X$, with $\ell=c_1(L)\in H^2(X,\mathbb Z)$.
Transport $\ell$ using the marking. A holomorphic line
bundle extending $L$ exists locally along precisely those
deformations for which $\ell$ stays of type $(1,1)$
(Lefschetz $(1,1)$ theorem).

\medskip
By Beauville--Bogomolov orthogonality this is
\[
\boxed{q(\ell,\sigma_t)=0.}
\]
Thus the locus of deformations of $(X,L)$ has period image
\[
\boxed{\Omega_{\Lambda,\ell}
=\Omega_\Lambda\cap\mathbb P(\ell^\perp).}
\]
For $\ell\ne0$ it is locally a smooth hypersurface of
dimension $b_2(X)-3=h^{1,1}(X)-1$.
\end{frame}

\begin{frame}{12. Polarized deformations: the geometric picture}
\small
\begin{center}
\begin{tabular}{c|l}
Object & What its points record\\ \hline
$\mathcal M_\Lambda$ & marked IHS manifolds $(X,\varphi)$\\
$\Omega_\Lambda$ & possible holomorphic symplectic lines\\
$\Omega_{\Lambda,\ell}$ & periods orthogonal to a fixed $\ell$
\end{tabular}
\end{center}
\medskip
The condition $q(\ell,\sigma)=0$ is a \emph{hyperplane}
condition in $\Lambda_{\mathbb C}$, but a
\emph{hypersurface} inside the period domain.

\medskip
If $L$ is ample, it stays ample under sufficiently small
deformations of the pair. Hence this hypersurface describes
the \emph{local moduli of polarized deformations} $(X_t,L_t)$.
For K3 surfaces: dimension $20$ drops to $19$.
\end{frame}

\end{document}
